\documentclass[10pt]{article}
\usepackage[UTF8]{ctex}
\usepackage{amsmath, amssymb}
\usepackage{geometry}
\geometry{a4paper, margin=1.5cm}
\usepackage{enumitem}

\title{导数知识清单——基础+应用完整版}
\author{}
\date{}

\begin{document}
\maketitle

\section*{致亲爱的同学}
本清单涵盖高中文科数学导数部分的核心知识，包括\textbf{基础概念与公式}和\textbf{应用方法}两大板块，按主题分级列出，便于复习和自查。愿你带着这份清单，稳扎稳打，将导数真正变为解题利器。

\section{基础知识}
\subsection{导数概念与定义}
\subsubsection{平均变化率}
函数 $y=f(x)$ 从 $x$ 到 $x+\Delta x$ 的平均变化率为
$$\displaystyle \frac{\Delta y}{\Delta x}=\frac{f(x+\Delta x)-f(x)}{\Delta x}.$$

\subsubsection{瞬时变化率——导数定义}
函数 $f(x)$ 在 $x$ 处的导数（瞬时变化率）：
$$\displaystyle f'(x)=\lim_{\Delta x\to 0}\frac{f(x+\Delta x)-f(x)}{\Delta x}.$$
若极限存在，则称 $f$ 在 $x$ 处可导。

\subsubsection{导数的几何意义}
导数 $f'(x_0)$ 表示曲线 $y=f(x)$ 在点 $(x_0,f(x_0))$ 处的\textbf{切线斜率}。

\subsection{基本初等函数的导数公式}
\begin{itemize}
    \item \textbf{常数函数}：$f(x)=c$（$c$ 为常数） $\to$ $f'(x)=0$
    \item \textbf{幂函数}：$f(x)=x^n$（$n\in\mathbb{Z}$） $\to$ $f'(x)=n x^{n-1}$；推广：对任意实数指数 $r$，$(x^r)'=r x^{r-1}$
    \item \textbf{指数函数}：$f(x)=e^x$ $\to$ $f'(x)=e^x$；$f(x)=a^x$（$a>0,a\neq1$） $\to$ $f'(x)=a^x\ln a$
    \item \textbf{对数函数}：$f(x)=\ln x$ $\to$ $f'(x)=\dfrac{1}{x}$；$f(x)=\log_a x$（$a>0,a\neq1$） $\to$ $f'(x)=\dfrac{1}{x\ln a}$
    \item \textbf{三角函数}：$f(x)=\sin x$ $\to$ $f'(x)=\cos x$；$f(x)=\cos x$ $\to$ $f'(x)=-\sin x$
\end{itemize}

\subsection{导数的四则运算法则}
设 $u(x),v(x)$ 可导，$c$ 为常数。
\begin{itemize}
    \item \textbf{和差法则}：$(u\pm v)' = u' \pm v'$
    \item \textbf{常数因子法则}：$(c u)' = c\, u'$
    \item \textbf{乘积法则}：$(uv)' = u'v + uv'$（口诀：前导后不导 + 后导前不导）
    \item \textbf{商法则}：$\left(\dfrac{u}{v}\right)' = \dfrac{u'v - uv'}{v^2}$（$v\neq0$）（口诀：分母平方，分子上导下不导减上不导下导）
\end{itemize}

\subsection{复合函数求导——链式法则}
若 $y=f(g(x))$，且 $u=g(x)$，则 $\dfrac{dy}{dx}=f'(g(x))\cdot g'(x)$。也可写作 $y'(x)=f'(u)\cdot u'(x)$。\textbf{关键}：由外到内逐层求导，各层导数相乘。

\subsection{常用极限与重要常数}
\begin{itemize}
    \item \textbf{自然常数} $e$：$e = \lim\limits_{n\to\infty}\left(1+\frac{1}{n}\right)^n \approx 2.71828$
    \item \textbf{重要极限}：
    $$\lim\limits_{x\to 0}\dfrac{\sin x}{x}=1,\quad \lim\limits_{x\to 0}\dfrac{\ln(1+x)}{x}=1\ \text{或}\ \lim\limits_{x\to 0}\dfrac{e^x-1}{x}=1.$$
\end{itemize}

\section{导数的应用}
\subsection{切线问题}
\subsubsection{已知切点求切线方程}
步骤：
\begin{enumerate}
    \item 求导 $f'(x)$；
    \item 计算斜率 $k=f'(x_0)$；
    \item 用点斜式 $y-f(x_0)=k(x-x_0)$ 写出方程。
\end{enumerate}
\begin{quote}
    \textbf{注意：} 先验证点 $(x_0,f(x_0))$ 是否在曲线上。
\end{quote}

\subsubsection{过曲线外一点求切线方程}
设切点法：
\begin{enumerate}
    \item 设切点 $(x_0,f(x_0))$；
    \item 写出切线方程 $y-f(x_0)=f'(x_0)(x-x_0)$；
    \item 将已知点代入，解出 $x_0$（可能多解）；
    \item 代回得切线方程。
\end{enumerate}
\begin{quote}
    \textbf{注意：} 可能有两条或更多切线，需全部求出。
\end{quote}

\subsection{单调性与极值}
\subsubsection{导数与单调性的关系}
在区间 $I$ 内：
\begin{itemize}
    \item 若 $f'(x)>0$，则 $f(x)$ 在 $I$ 上\textbf{单调递增}；
    \item 若 $f'(x)<0$，则 $f(x)$ 在 $I$ 上\textbf{单调递减}；
    \item 若 $f'(x)=0$，可能是极值点或平台。
\end{itemize}

\subsubsection{求单调区间的步骤}
\begin{enumerate}
    \item 确定定义域；
    \item 求导 $f'(x)$；
    \item 解 $f'(x)=0$，得临界点；
    \item 列表，判断各区间内 $f'(x)$ 的符号；
    \item 写出单调区间（通常用开区间）。
\end{enumerate}

\subsubsection{极值的判定}
\begin{itemize}
    \item \textbf{必要条件}：极值点处 $f'(x)=0$ 或导数不存在。
    \item \textbf{充分条件（一阶导数法）}：
    \begin{itemize}
        \item 若 $x_0$ 左侧 $f'(x)>0$，右侧 $f'(x)<0$，则 $x_0$ 是\textbf{极大值点}；
        \item 若左侧 $f'(x)<0$，右侧 $f'(x)>0$，则 $x_0$ 是\textbf{极小值点}。
    \end{itemize}
\end{itemize}

\subsubsection{求极值的步骤}
\begin{enumerate}
    \item 求导，找临界点；
    \item 列表判断每个临界点两侧导数符号；
    \item 得出极值点及极值。
\end{enumerate}

\subsection{最值问题}
\subsubsection{闭区间上的最值}
步骤：
\begin{enumerate}
    \item 求导，找出区间内的临界点；
    \item 计算区间端点函数值和各临界点函数值；
    \item 比较得最大值和最小值。
\end{enumerate}
\begin{quote}
    \textbf{注意：} 最值可能出现在端点，不可遗漏。
\end{quote}

\subsubsection{实际优化问题}
流程：
\begin{enumerate}
    \item 设变量（通常一个自变量）；
    \item 根据几何或实际关系建立目标函数；
    \item 确定自变量的实际范围（定义域）；
    \item 求导，找临界点；
    \item 比较端点与临界点，得最值；
    \item 回答实际问题（带单位）。
\end{enumerate}
常见模型：面积最大、容积最大、成本最低、利润最大等。

\subsection{含参讨论}
\subsubsection{含参单调性讨论}
若 $f'(x)$ 是二次函数，需根据判别式讨论 $f'(x)\ge 0$ 恒成立的条件。典型问题：已知函数在 $\mathbb{R}$ 上单调递增，求参数范围。

\subsubsection{含参极值讨论}
当 $f'(x)=0$ 的解含参数时，需根据参数范围讨论极值点的个数和性质。分类标准：二次项系数、判别式、根的大小关系等。

\subsubsection{分类讨论的书写规范}
\begin{itemize}
    \item 明确分类标准（如参数的正负、与临界值的大小关系）；
    \item 每种情况独立判断；
    \item 最后归纳结论。
\end{itemize}

\section*{使用建议}
\begin{itemize}
    \item \textbf{基础部分}：务必熟练掌握，这是求导的工具箱。
    \item \textbf{应用部分}：对应高考核心考点，需结合典型例题反复练习。
    \item \textbf{复习时}：先对照清单回顾知识点，再通过综合练习查漏补缺。
\end{itemize}

\vspace{1cm}
\begin{center}
    \large 祝你学习顺利，导数不再困难！
\end{center}

\end{document}